\documentclass[11pt,a4paper]{article}

\usepackage{fontspec}
\setmainfont{DejaVu Sans}
\setsansfont{DejaVu Sans}
\setmonofont{DejaVu Sans Mono}[Scale=0.88]
\defaultfontfeatures{Ligatures=TeX}
\usepackage{microtype}
\usepackage{unicode-math}
\setmathfont{Latin Modern Math}[Scale=1.0]
\usepackage[ngerman]{babel}
\usepackage{geometry}
\geometry{margin=2.5cm}

% --- Zeilenumbruch: sauber wie im Harmonikabuch ---
\sloppy
\emergencystretch=3em
\tolerance=2500
\hbadness=10000
\hfuzz=2pt
\clubpenalty=10000
\widowpenalty=10000

\usepackage{amsmath,amssymb}
\usepackage{parskip}
\usepackage{booktabs}
\usepackage{array}
\usepackage{xcolor}
\usepackage{hyperref}
\hypersetup{
  colorlinks=true,
  linkcolor=blue!50!black,
  urlcolor=blue!50!black,
  pdftitle={Das primordiale Lithium-Problem und FFGFT},
  pdfauthor={Johann Pascher}
}

\title{Das primordiale Lithium-Problem und FFGFT\\[2pt]
  \large Kurzfassung für die IPI-Korrespondenz}
\author{Johann Pascher\\[2pt]
\small ORCID: \href{https://orcid.org/0009-0000-6518-4064}{0009-0000-6518-4064}\\
\small GitHub: \url{https://github.com/jpascher/T0-Time-Mass-Duality}\\
\small Zenodo (FFGFT v1.0.13): \url{https://zenodo.org/records/20041529}\\
\small DOI: \href{https://doi.org/10.5281/zenodo.20041529}{10.5281/zenodo.20041529}}
\date{Mai 2026}

\begin{document}
\maketitle

\begin{abstract}
\noindent Das primordiale Lithium-Problem -- die hartnäckige Faktor-3--4-Diskrepanz ($4{,}2\sigma$--$5{,}3\sigma$) zwischen den Vorhersagen der Big-Bang-Nukleosynthese (BBN) für $^7$Li und Beobachtungen metallarmer Halo-Sterne (das Spite-Plateau) -- hat sich seit vier Jahrzehnten der Auflösung widersetzt. Die Standard-Lösungsklassen (astrophysikalische Depletion, Kernphysik, Physik jenseits des Standardmodells) stoßen alle auf strukturelle Schwierigkeiten. Fields (2011) listet ausdrücklich ``große Änderungen am kosmologischen Rahmen'' als mögliche Lösungsklasse auf. FFGFT instanziiert diese Klasse: In einem statischen, ewigen $\xi$-Universum gibt es keine primordiale Nukleosynthese-Epoche, und das Spite-Plateau ist nicht das Relikt eines Expansionsuniversum-Ereignisses. Diese Notiz präsentiert die FFGFT-Position einschließlich einer vorläufigen Erklärungs-Skizze, warum die Häufigkeiten von D und $^4$He dennoch mit BBN+CMB-Vorhersagen übereinstimmen, $^7$Li jedoch nicht. Die Erklärung wird als Skizze, nicht als abgeschlossene Herleitung präsentiert.
\end{abstract}

\section{Das Problem in Zahlen}

Die Big-Bang-Nukleosynthese-Vorhersage liefert mit dem WMAP/Planck-Baryon-zu-Photon-Verhältnis $\eta_{10} = 6{,}23 \pm 0{,}17$:
\begin{align}
\text{D/H}\big|_{\text{BBN+CMB}} &= (2{,}82 \pm 0{,}21) \times 10^{-5} &&\text{(Beobachtung: passt)}\\
Y_p(^4\text{He}) &= 0{,}249 \pm 0{,}009 &&\text{(Beobachtung: passt)}\\
^{7}\text{Li/H}\big|_{\text{BBN+CMB}} &= (5{,}24^{+0{,}71}_{-0{,}67}) \times 10^{-10} &&\text{(Beobachtung: weicht ab)}
\end{align}
Die Spite-Plateau-Messung \cite{spitespite1982,fields2011} metallarmer Halo-Sterne ergibt
\begin{equation}
^{7}\text{Li/H}\big|_{\text{obs}} = (1{,}23^{+0{,}68}_{-0{,}32}) \times 10^{-10}\,,
\end{equation}
einen Faktor von $2{,}4$--$4{,}3$ unterhalb der BBN+CMB-Vorhersage, entsprechend $4{,}2\sigma$--$5{,}3\sigma$ \cite{cyburt2008}. Cyburt, Fields und Olive betitelten ihr Update von 2008 ``Eine bittere Pille: Das primordiale Lithium-Problem verschlimmert sich'' -- verbesserte Daten zu nuklearen Wirkungsquerschnitten haben die vorhergesagte $^7$Li-Häufigkeit \emph{erhöht}.

Die Diskrepanz ist universell über stellare Umgebungen hinweg: Das Spite-Plateau wurde in Halo-Feldsternen, Kugelhaufen \cite{cyburt2008}, Zwergsphäroidalen, ultra-faint-Galaxien und akkretierten Satelliten \cite{borisov2024} beobachtet. Es ist kein Merkmal einer einzelnen galaktischen Umgebung.

\section{Standard-Lösungsklassen}

Nach Fields (2011) wurden drei Klassen verfolgt:

\paragraph{Astrophysikalisch (stellare Depletion).}
Das aktivste aktuelle Programm \cite{borisov2024} modelliert interne stellare Mischung -- atomare Diffusion, parametrische Turbulenz, rotationsinduzierten Drehimpulstransport -- um photosphärisches $^7$Li um einen Faktor $\sim 3$ zu reduzieren. Die Herausforderung: Das Spite-Plateau ist beobachterisch \emph{dünn}, mit wenig Streuung von Stern zu Stern, und es findet sich über sehr verschiedene Sternpopulationen hinweg. Universelle Depletion über diese Umgebungen ist physikalisch schwer zu motivieren.

\paragraph{Kernphysik.}
Hypothetische Resonanzen, die die $^7$Be-Zerstörung verstärken, z.\,B.\ $^7$Be$ + ^{3}$He $\to {}^{10}$C$^* \to p + {}^{9}$B \cite{fields2011,hammache2013}, oder revidierte $^3$He$(\alpha,\gamma)^7$Be-Wirkungsquerschnitte. Jüngste experimentelle Arbeit hat die Raten \emph{verschärft} und die Vorhersage \emph{verschlechtert} \cite{cyburt2008}.

\paragraph{Jenseits des Standardmodells.}
Zerfallende SUSY-Teilchen, zeitlich variable Naturkonstanten \cite{uzan2003}, primordiale Magnetfelder. Constraint: D und $^4$He dürfen nicht gestört werden, was ein enges Parameterfenster lässt.

\paragraph{Die vierte Klasse.}
Fields (2011) schließt seine Klassifikation mit: \emph{``große Änderungen am kosmologischen Rahmen, der zur Interpretation der Daten leichter Elemente verwendet wird.''} Diese Kategorie hat weniger Beachtung erhalten, ist aber logisch offen.

\section{Die FFGFT-Position}

FFGFT instanziiert die vierte Klasse. Der kosmologische Rahmen ist statisch und ewig, abgeleitet aus dem einzigen dimensionslosen Parameter $\xi = 4/30{,}000$:
\begin{itemize}\itemsep1pt
\item Kein Urknall. Das Heisenberg-Zeit-Energie-Unschärfe-Argument (Doc 025) schließt eine kosmologische Singularität endlicher Dauer aus.
\item Keine expandierende Raumzeit. Die kosmologische Rotverschiebung ist ein geometrischer Pfad-Effekt durch das körnige $\xi$-Vakuum, keine Doppler-Signatur (Doc 026).
\item Keine primordiale Nukleosynthese-Epoche. Das Spite-Plateau ist daher \emph{nicht} das Relikt eines universumsweiten Ereignisses vor 13,8 Mrd.\ Jahren.
\end{itemize}

In diesem Rahmen verzeichnet das Spite-Plateau die langfristige Gleichgewichts-$^7$Li-Häufigkeit, die durch stellare Prozesse (galaktische Kosmische-Strahlen-Spallation, Novae, AGB-Sterne, Kernkollaps-Supernovae, Weißer-Zwerg-Mischen \cite{borisov2024}) in einem quasi-stationären Zustand produziert und aufrechterhalten wird. Die Faktor-3--4-``Diskrepanz'' ist dann keine Diskrepanz mit der Realität, sondern mit einer Vorhersage, die ein Ereignis (BBN) erfordert, das in diesem Rahmen nicht stattfand.

Die ehrliche Folgefrage lautet: \emph{Warum stimmen D und $^4$He so präzise mit BBN+CMB überein, wenn es keine BBN gibt?}

\section{Erklärungs-Skizze: Warum D und \texorpdfstring{$^4$He}{4He} übereinstimmen}

Was folgt, ist eine \emph{Erklärungs-Skizze}, keine abgeschlossene Herleitung. Der Punkt besteht darin, einen plausiblen Mechanismus anzudeuten, nicht zu behaupten, der Mechanismus sei bereits ausgearbeitet.

\paragraph{Schritt 1 -- Die CMB ist real, aber kein Relikt.}
FFGFT akzeptiert die CMB als physikalisches Phänomen, identifiziert sie aber als Manifestation des $\xi$-Vakuums statt als rotverschobenes Nachglühen einer heißen Frühphase. In Doc~091 wird die CMB-Energiedichte festgelegt durch
\begin{equation}
\rho_{\mathrm{CMB}} = \frac{\xi\,\hbar c}{L_\xi^4}\,, \qquad L_\xi \approx 100{,}24\,\mu\text{m}\,,
\end{equation}
und die CMB-Temperatur folgt entsprechend. Das Photonen-Bad, das das Spektrum trägt, ist real und operativ \emph{jetzt}, nicht nur bei $z \approx 1100$.

\paragraph{Schritt 2 -- Das Baryon-zu-Photon-Verhältnis.}
In der Standard-BBN hängen die Häufigkeiten leichter Elemente vom einzigen Parameter $\eta = n_B/n_\gamma$ ab. In FFGFT sind sowohl die Photonendichte (aus $\rho_{\mathrm{CMB}}$ und $T_{\mathrm{CMB}}$) als auch die kosmische Baryonendichte (aus kontinuierlichen stellaren Prozessen, integriert über ein ewiges Universum) gegenwärtige Größen. Ihr Verhältnis definiert ein \emph{effektives} $\eta_{\mathrm{eff}}$, das die WMAP/Planck-Inferenz korrekt aus den CMB-akustischen Peaks extrahiert -- weil die akustische-Peak-Struktur nur vom gegenwärtigen $\eta$ abhängt, nicht von dessen Geschichte.

Die ΛCDM-Interpretation rechnet von diesem $\eta$ über angenommene Entropie-Erhaltung auf ein primordiales $\eta$ zurück; FFGFT behält das gegenwärtige $\eta$ als stationäres Charakteristikum des $\xi$-Universums.

\paragraph{Schritt 3 -- Thermische Gleichgewichts-Häufigkeiten.}
Für leichte Nuklide, deren Häufigkeiten durch \emph{schnelle Zwei-Körper-Gleichgewichtsreaktionen in einem thermischen Bad} bestimmt werden, ist die Gleichgewichts-Häufigkeit nur eine Funktion von $\eta$ und $T$, nicht der Geschichte, die zu diesem $(\eta, T)$-Paar geführt hat. Das Reaktionsnetzwerk, das D- und $^4$He-Häufigkeiten unter BBN-Bedingungen bestimmt, wird von solchen Zwei-Körper-Gleichgewichten mit starken Rückreaktionen dominiert:
\begin{equation}
p + n \leftrightarrow d + \gamma, \qquad d + d \leftrightarrow {}^3\text{He} + n, \qquad {}^3\text{He} + d \leftrightarrow {}^4\text{He} + p\,, \dots
\end{equation}
Wenn das FFGFT-$\xi$-Vakuum ein $(T_{\mathrm{CMB}}, \eta_{\mathrm{eff}})$ aufrechterhält, das mit den WMAP-inferierten Werten übereinstimmt, folgen dieselben Gleichgewichts-Häufigkeiten. Daher stimmen die Beobachtungen von D und $^4$He mit BBN+CMB-Vorhersagen überein: Die Vorhersagen beschreiben ein \emph{thermisches Gleichgewicht}, und das Gleichgewicht wird in FFGFT für dasselbe $(T, \eta)$ konstruktionsbedingt realisiert. Die Übereinstimmung ist kein Zufall; sie ist ein Konsistenztest der $\xi$-Bestimmung von $T_{\mathrm{CMB}}$.

\paragraph{Schritt 4 -- Warum $^7$Li das Muster bricht.}
$^7$Li ist die Ausnahme, weil seine Produktions-Zerstörungs-Kette einen \emph{langsamen, Nicht-Gleichgewichts}-Anteil aufweist:
\begin{enumerate}\itemsep1pt
\item $^7$Li wird bei WMAP-$\eta$ überwiegend als $^7$Be via $^3$He$(\alpha,\gamma)^7$Be produziert \cite{fields2011}.
\item In der Standard-BBN überlebt $^7$Be dann das Erlöschen der Kernreaktionen und zerfällt durch Elektronen-Einfang über eine Zeitskala von $\sim 53$ Tagen zu $^7$Li -- nach dem Freeze-out der BBN.
\item In FFGFT findet kein Freeze-out statt, weil es keinen expansionsbedingten Temperaturabfall gibt. Die $^7$Be-Population sitzt in einer stationären Umgebung mit kontinuierlichen Photon- und Elektronen-Wechselwirkungen, einschließlich Kanälen, die $^7$Be auf Zeitskalen viel kürzer zerstören, als seine Labor-Halbwertszeit nahelegen würde.
\end{enumerate}
Die BBN-Berechnung nimmt eine saubere $^7$Be $\to {}^7$Li-Konversion nach Freeze-out an und behandelt sie als Buchführung. In FFGFT ist diese Annahme ungültig: Der stationäre Zustand erlaubt kontinuierliche $^7$Be-Zerstörung (über kosmische-Strahlen-induzierte Spallation, energetische Photonen-Kanäle im $\xi$-Vakuum oder stellare Inkorporation), die das BBN-Bookkeeping nicht berücksichtigt. Die Plateau-Niveau-$^7$Li-Häufigkeit ist dann der stationäre Attraktor des vollen Netzwerks in einem statischen Universum, und sie ist natürlicherweise niedriger als der BBN-Bookkeeping-Wert.

\paragraph{Was diese Skizze beansprucht und nicht beansprucht.}
\textit{Beansprucht:} ein struktureller Grund, warum ein gleichgewichtsdominiertes Nuklid (D, $^4$He) der Standard-Vorhersage folgt, während ein freeze-out-abhängiges Nuklid ($^7$Li) dies nicht tut. \textit{Nicht beansprucht:} eine geschlossene quantitative Herleitung des Spite-Plateau-Werts aus $\xi$. Der nächste Schritt wäre, das stationäre Netzwerk im $\xi$-Vakuum aufzustellen und den Attraktor für $^7$Li zu berechnen. Dies ist offene Arbeit, kein abgeschlossenes Ergebnis.

\section{Zusammenfassung}

\begin{enumerate}\itemsep2pt
\item Das Lithium-Problem ist real und etabliert: Faktor 3--4, $4{,}2\sigma$--$5{,}3\sigma$, hartnäckig über vier Jahrzehnte und viele Lösungsversuche.
\item Astrophysikalische, kernphysikalische und BSM-Lösungen stoßen alle auf strukturelle Herausforderungen; die Alternative über den kosmologischen Rahmen (Fields 2011) bleibt offen.
\item FFGFT instanziiert die Alternative über den kosmologischen Rahmen: kein Urknall, keine primordiale Nukleosynthese-Epoche, das Spite-Plateau ist ein stationäres stellares Gleichgewichts-Merkmal.
\item Die Übereinstimmung von D und $^4$He mit BBN+CMB ist nicht zufällig, sondern ein Konsistenztest: Gleichgewichts-Häufigkeiten hängen nur von $(T, \eta)$ ab, und FFGFT legt beides über die $\xi$-Geometrie fest.
\item $^7$Li weicht ab, weil seine Häufigkeit von einer Freeze-out-Annahme abhängt, die in einem statischen Universum ungültig ist; der stationäre Attraktor liegt unterhalb des BBN-Bookkeeping-Werts.
\item Die quantitative Herleitung des FFGFT-$^7$Li-Attraktors ist offene Arbeit; diese Notiz präsentiert die Erklärung als Skizze.
\end{enumerate}

\section*{Referenzen}

\noindent\small
\begin{thebibliography}{99}
\bibitem{spitespite1982} Spite, M.\ \& Spite, F.\ (1982). \emph{Abundance of lithium in unevolved halo stars and old disk stars.} A\&A 115, 357.
\bibitem{nollett1997} Nollett, K.M., Lemoine, M.\ \& Schramm, D.N.\ (1997). \emph{Big Bang Nucleosynthesis in the new cosmology.} arXiv:astro-ph/9612197.
\bibitem{dring1997} Möring \emph{et al.} (1997). \emph{Reaction rate measurements relevant to BBN $^7$Li production.} ApJ.
\bibitem{cyburt2008} Cyburt, R.H., Fields, B.D.\ \& Olive, K.A.\ (2008). \emph{A Bitter Pill: The Primordial Lithium Problem Worsens.} arXiv:0808.2818.
\bibitem{dunkley2009} Dunkley, J.\ \emph{et al.} (2009). \emph{Five-Year WMAP Observations: Bayesian Estimation of CMB Polarization Maps.} arXiv:0803.0586.
\bibitem{fields2011} Fields, B.D.\ (2011). \emph{The Primordial Lithium Problem.} Annual Reviews of Nuclear and Particle Science. arXiv:1203.3551.
\bibitem{hammache2013} Hammache, F.\ \emph{et al.} (2013). \emph{High-energy direct measurement of the $^7$Be$(p,\gamma)^8$B reaction at LUNA.} arXiv:1312.0894.
\bibitem{uzan2003} Uzan, J.-P.\ (2003). \emph{The fundamental constants and their variation: observational status and theoretical motivations.} arXiv:hep-ph/0205340.
\bibitem{borisov2024} Borisov, S., Charbonnel, C., Prantzos, N., Dumont, T.\ \& Palacios, A.\ (2024). \emph{A coherent view of Li depletion and angular momentum transport to explain the Li plateau.} A\&A. arXiv:2403.15534.
\end{thebibliography}

\medskip
\noindent\emph{FFGFT-Quelldokumente: Doc~025 (T0-Kosmologie, $\xi$-Feld-Universum und CMB), Doc~026 (geometrische Kosmologie, $H_0$ und der Rotverschiebungs-Mechanismus), Doc~091 (Casimir-CMB-Beziehung, $L_\xi$). Alle verfügbar unter \url{https://github.com/jpascher/T0-Time-Mass-Duality} und Zenodo v1.0.13: \url{https://zenodo.org/records/20041529}.}

\end{document}
